\chapter{Tools, testing, and on-device export}

A research codebase is more than the model. This chapter surveys the supporting cast:
the runnable scripts, the test suite, checkpoint migration, and the export path that puts
the model on an embedded device.

\section{The scripts you run}
All under \file{tools/} and \file{scripts/}; full details and copy-paste commands are in
\file{docs/scripts.md}.

\begin{center}
\renewcommand{\arraystretch}{1.2}
\begin{tabular}{l >{\raggedright\arraybackslash}p{8.8cm}}
\toprule
\textbf{Command} & \textbf{What it does} \\
\midrule
\code{scripts.run\_train} & launch training (prefer \code{tools/train\_supervisor.sh} for long runs) \\
\code{tools.eval} & COCO / polygon / distance metrics on a checkpoint \\
\code{tools.infer} & draw box + polygon overlays on arbitrary images \\
\code{tools.continuous\_eval} & auto-evaluate new checkpoints in a run directory \\
\code{tools.export\_saved\_model} & host/server SavedModel (NMS baked in, $[0,1]$ input) \\
\code{tools.device.export\_device\_dlc} & on-device SNPE/DLC export \\
\code{tools.checkpoint\_migration} & migrate / warm-start checkpoints \\
\code{tools.benchmark\_pipeline} & data-pipeline throughput \\
\code{tools.pipeline.diagnose\_pipeline} & stage-by-stage pipeline timing \\
\bottomrule
\end{tabular}
\end{center}

\section{Checkpoint migration and warm-starting}
\file{tools/checkpoint\_migration.py} loads weights into the model from an existing
checkpoint, auto-detecting the kind:

\begin{center}
\renewcommand{\arraystretch}{1.2}
\begin{tabular}{l l >{\raggedright\arraybackslash}p{5.0cm}}
\toprule
\textbf{Source} & \textbf{Strategy} & \textbf{What transfers} \\
\midrule
legacy TF2-Vision ckpt & \code{frozen} & backbone + decoder (+ head if classes match) \\
this codebase's own ckpt & \code{native} & requested modules, from EMA weights \\
\bottomrule
\end{tabular}
\end{center}

\begin{pitfall}
Warm-start from a checkpoint that carries EMA (a periodic \code{ckpt-N} or a \code{best\_*}),
\textbf{not} a bare model export. The plain model object graph omits some list-tracked C2f
block variables (a Keras quirk), so only the EMA shadows hold the \emph{complete} weights.
The \code{native} strategy loads via that EMA path for exactly this reason.
\end{pitfall}

\section{Testing}
The \code{pytest} suite (\file{tests/}) runs in \textbf{eager mode} and is split into
\code{unit/} (components --- backbone, decoders, EMA, the assigner, the evaluators),
\code{integration/} (full pipeline, checkpoint migration, a real 2-replica multi-GPU test
on virtual CPU devices), \code{smoke/} (a 10-step training loop), and top-level component
tests (losses, mosaic, parser, polygon preprocessing).

\begin{hood}
Two conventions make the tests trustworthy. (1) Unit/integration tests \textbf{never}
require the real datasets --- they build synthetic tensors inline, so CI runs anywhere.
(2) Tests are \emph{discriminating}: they pin a value or relationship that \emph{fails} on
the buggy behavior and passes once fixed (see \file{tests/test\_loss\_reference\_parity.py}),
rather than merely asserting ``runs without error''.
\end{hood}

\begin{lstlisting}[language=bash]
# fast, no datasets (what CI runs):
pytest tests/unit tests/integration -q
# the multi-GPU test needs a fresh process:
pytest tests/integration/test_multigpu.py -q
\end{lstlisting}

\section{On-device export (Qualcomm SNPE / DLC)}
The model ultimately runs on an embedded device. \file{tools/device/export\_device\_dlc.py}
emits a TensorFlow SavedModel laid out as a \textbf{drop-in replacement} for a legacy
on-device DLC, so the existing SNPE conversion pipeline keeps working unchanged. It differs
from the host export in deliberate ways: it bakes the $/255$ normalization in (the device
feeds raw $[0,255]$ pixels), runs in \code{float32}, emits one raw tensor per head (FPN
levels concatenated), and promotes each head to a cleanly named top-level graph node the
SNPE extractor expects.

\begin{pitfall}
The deployed device decoder stores anchors as $(y,x)$ and expects box channels in
\textbf{y-first} order $[t,l,b,r]$, while the model (and the host code) is x-first
$[l,t,r,b]$. The exporter therefore \textbf{reorders the box channels} by default
(\code{-{}-legacy\_box\_order}). Feeding x-first boxes to the y-first decoder transposes
every box --- the kind of bug that drops on-device accuracy from $0.68$ to $0.19$ while
every shape still looks correct. The exporter's \code{-{}-verify} flag checks this and more.
\end{pitfall}
